% Part: applied-modal-logics
% Chapter: epistemic-logic
% Section: relational-models

\documentclass[../../../include/open-logic-section]{subfiles}

\begin{document}

\olfileid{aml}{el}{rel}

\olsection{Model Relasional}

Gagasan semantik dhasar kanggo logika epistemik padha karo logika
modal lumrah. Model relasional tetep ngemot himpunan jagad lan
panetepan kang nemtokake !!{propositional
variable}s endi kang
dianggep ``bener'' ing saben jagad. Yen mung ana siji agen, kita
nduweni siji relasi aksesibilitas kaya biyasane. Nanging, ing
logika epistemik multiagen, relasi aksesibilitas siji mau diganti
himpunan relasi aksesibilitas, siji kanggo saben~$a$ ing himpunan
lambang agen~$G$.

\emph{Model relasional} ngemot himpunan jagad kang disambungake
dening relasi aksesibilitas biner---siji kanggo saben agen---lan
panetepan kang nemtokake !!{propositional variable}s endi kang
bener ing saben jagad.

\begin{defn}
  \emph{Model} kanggo basa epistemik multiagen yaiku tripel
  $\mModel{M} = \tuple{W, R, V}$, where
  \begin{enumerate}
  \item $W$ iku himpunan ``jagad'' kang ora kosong,
  \item kanggo saben $a \in G$, ${R}_a$ iku relasi aksesibilitas biner
  ing~$W$, lan
  \item $V$ iku fungsi kang netepake marang saben !!{propositional
    variable}~$p$ sawijining himpunan jagad mungkin $V(p)$.
  \end{enumerate}
  Yen $R_a ww'$, kita kandha yen $w'$ \emph{bisa diakses dening a
    saka}~$w$. Yen $w \in V(p)$, kita kandha $p$ \emph{bener ing}~$w$.
\end{defn}

Cara kerjane padha karo logika modal normal, mung relasi
aksesibilitase luwih akeh. Kanggo sawijining agen, relasi
aksesibilitase lumrahe ditafsirake minangka gegambaran kahanan
informasi agen iku. Contone, $R_a ww'$ kerep dianggep medharake
yen $w'$ selaras karo informasi kang diduweni~$a$ ing~$w$.
Kanthi tembung liya, ing~$w$, agen iku ora bisa mbedakake jagad
$w$ lan jagad~$w'$.

\end{document}
